\textbf{Equação de crescimento composto:}

Após $n$ anos com taxa anual $r$, um capital $C_0$ torna-se:
\[
C(n) = C_0(1+r)^n.
\]

Para dobrar em $n=7$ anos:
\[
2\,C_0 = C_0(1+r)^7 \implies (1+r)^7 = 2.
\]

\textbf{Resolvendo para $r$:}
\[
1+r = 2^{1/7} = e^{\ln 2/7}.
\]

Calculando $2^{1/7}$:
\[
2^{1/7} = e^{\ln 2/7} = e^{0{,}6931/7} = e^{0{,}09902} \approx 1 + 0{,}09902 + \frac{(0{,}09902)^2}{2} + \cdots \approx 1{,}09902.
\]

Portanto:
\[
r = 2^{1/7}-1 \approx 0{,}09902 \approx \mathbf{9{,}90\%\ \text{ao ano}}.
\]

\textbf{Verificação:}
\[
(1{,}099)^7 = e^{7\ln(1{,}099)} = e^{7\times0{,}09441} = e^{0{,}6609}\approx1{,}936\approx2.\quad\checkmark
\]

\textbf{Regra dos 72 (aproximação prática):}

A taxa aproximada para dobrar em $n$ anos é $r\approx72/n$:
\[
r\approx\frac{72}{7}\approx10{,}3\%.
\]
Essa regra superestima levemente o valor exato de $9{,}90\%$.